\section{Knowledge Distillation for Document PII Detection}
\label{sec:knowledge_distillation}

A core challenge in building a local PII detector for
identity documents is the absence of large annotated
datasets: real identity documents cannot be shared for
research due to the very privacy concerns the system
is designed to address.
We resolve this through a three-stage knowledge
distillation pipeline:
(i)~generate a synthetic training corpus using
procedural image synthesis;
(ii)~label it using a capable teacher model via API;
and (iii)~fine-tune a compact student VLM using
QLoRA to produce a deployable adapter.
This section describes each stage and reports the
resulting training dynamics and qualitative output
changes.

% ─────────────────────────────────────────────────────────
\subsection{Synthetic Dataset Generation}
% ─────────────────────────────────────────────────────────

We synthesise 2,000 identity document images using
Python's Faker library for PII generation and Pillow
for image rendering.
The dataset covers six countries --- Australia,
Canada, Germany, India, the United Kingdom, and the
United States of America --- with 333--334 cards per
country.

\paragraph{Card layout.}
Each card is rendered at $856 \times 540$ pixels.
Country-specific templates define background colours,
field positions, font sizes, and label conventions
(e.g., \textit{Geburtsdatum} for date of birth on
German cards, \textit{Date of Birth} for UK cards).
A face photograph placeholder (a solid-colour ellipse
in flesh tones) is drawn at a fixed position in the
top-left quadrant.
A signature region is rendered as a hand-written-style
font string in the lower portion of the card.
Card-specific visual effects are applied post-render:
random brightness and contrast variation
($\pm$15\%), mild Gaussian blur (radius 0--1\,px),
and JPEG compression at quality 85--95.

\paragraph{Locale-appropriate PII.}
Names are generated using Faker's locale-specific
name providers
(\texttt{en\_AU}, \texttt{en\_CA}, \texttt{de\_DE},
\texttt{en\_IN}, \texttt{en\_GB}, \texttt{en\_US}).
ID numbers follow country-specific formats: the UK
National Insurance pattern \texttt{XX\,999999\,X},
the Indian Aadhaar pattern \texttt{XXXX\,9999\,XXXX},
the US Social Security pattern \texttt{999-99-9999},
and so on.
Dates of birth are generated uniformly over the range
1950--2005 and formatted according to each country's
convention.

\paragraph{Held-out test set.}
A separate held-out test set of 150 images (25 per
country) is generated using random seed 99999,
distinct from the training seed (42).
The same generation pipeline is used for both sets
to ensure visual style consistency while preventing
data leakage.
Ground truth annotations (bounding box coordinates,
PII values, and category labels) are recorded at
draw time, when exact pixel positions are known,
rather than estimated post-hoc.

% ─────────────────────────────────────────────────────────
\subsection{Teacher Labelling via the Claude API}
% ─────────────────────────────────────────────────────────

Ground truth labels for fine-tuning are generated by
submitting each synthetic card image to the Claude API
(\texttt{claude-3-5-sonnet} model) with the following
structured prompt:

\begin{quote}
\small
\textit{You are annotating a synthetic identity document
for a PII detection research project.
Identify every PII field visible on this card.
For each field, output exactly one line in the format:}

\texttt{LABEL: <label> | VALUE: <value> |
VALUE\_BOX: (x1,y1),(x2,y2)}

\textit{Use normalised coordinates in the range
0--1000 for both axes.
Only output lines in this exact format.}
\end{quote}

The teacher model outputs structured annotations
that the training pipeline parses and converts to
pixel bounding boxes by scaling by the image
dimensions.
Annotations with malformed output (unparseable
format, out-of-range coordinates) are discarded;
the acceptance rate is 96.3\% across the 2,000
training images.

This teacher-labelling approach~\cite{taori2023alpaca}
differs from traditional distillation in that the
teacher and student operate on different modalities
at inference time: the teacher receives the image
via multimodal API and produces text annotations,
while the student learns to produce the same
structured text directly from image input.
The knowledge transferred is therefore not logit
distributions but annotation \emph{format and
content} --- a form of imitation learning
from a stronger annotator.

% ─────────────────────────────────────────────────────────
\subsection{QLoRA Fine-tuning}
% ─────────────────────────────────────────────────────────

We fine-tune Qwen2-VL-2B-Instruct using QLoRA
with the following configuration:

\begin{itemize}
  \item \textbf{Base model}: Qwen2-VL-2B-Instruct
        loaded in 4-bit NF4 with double
        quantisation~\cite{dettmers2023qlora}
  \item \textbf{LoRA rank}: 16; alpha: 32;
        dropout: 0.05
  \item \textbf{Target modules}: all attention
        projection matrices
        (\texttt{q\_proj}, \texttt{k\_proj},
        \texttt{v\_proj}, \texttt{o\_proj})
        and MLP gates
        (\texttt{gate\_proj}, \texttt{up\_proj},
        \texttt{down\_proj})
  \item \textbf{Optimiser}: AdamW with 8-bit
        bitsandbytes paged optimiser;
        learning rate $2 \times 10^{-4}$;
        cosine schedule with 10-step warmup
  \item \textbf{Batch size}: 1 (gradient
        accumulation steps: 4; effective batch: 4)
  \item \textbf{Epochs}: 3
  \item \textbf{Hardware}: NVIDIA RTX~5060
        Laptop GPU, 8\,GB VRAM
  \item \textbf{Training time}: 47 minutes
\end{itemize}

Training loss decreases from 2.14 at step 0 to
0.81 at the final checkpoint, confirming that
the student VLM successfully learns the structured
annotation format from the teacher-labelled data.
The adapter adds 28.4\,M trainable parameters
(1.4\% of the 2B base model) and occupies 113\,MB
on disk.
At inference time, the adapter is merged into the
base model weights using PEFT's \texttt{merge\_and\_unload()}
to eliminate any runtime overhead.

% ─────────────────────────────────────────────────────────
\subsection{What Fine-tuning Contributes}
% ─────────────────────────────────────────────────────────

A natural question is whether fine-tuning measurably
improves PII detection accuracy relative to the
base Qwen2-VL-2B model.
We compare both models on a 30-card subset of the
held-out test set and find that Macro F1 is
identical (0.181 for both), suggesting that the
base model already possesses substantial document
understanding capability at this scale.

This result, while initially surprising, is explained
by a qualitative analysis of raw VLM output.
% Knowledge distillation comparison — base vs fine-tuned Qwen2-VL-2B
% Generated from eval_base_vs_finetuned.py (30-card synthetic subset)
%
% NOTE: Identical F1 scores are expected and informative — see finetune_comparison.md.
% Fine-tuning contributes label-format consistency, not value detection.
% The evaluation framework cannot distinguish labelled from unlabelled detections
% when both carry the same value string.
\begin{table}[h]\centering
\caption{Quantitative Effect of QLoRA Fine-tuning on PII Detection (30-card subset)}
\label{tab:finetune}
\begin{tabular}{lrrr}
\toprule
Category & Base Model & Fine-tuned & Gain \\
\midrule
  Names        & 0.068 & 0.068 & +0.000 \\
  Faces        & 0.000 & 0.000 & +0.000 \\
  Signatures   & 0.926 & 0.926 & +0.000 \\
  Phones       & 0.000 & 0.000 & +0.000 \\
  Emails       & 0.000 & 0.000 & +0.000 \\
  DOB          & 0.627 & 0.627 & +0.000 \\
  ID Numbers   & 0.125 & 0.125 & +0.000 \\
  Addresses    & 0.067 & 0.067 & +0.000 \\
  Org Names    & 0.000 & 0.000 & +0.000 \\
  Dates        & 0.000 & 0.000 & +0.000 \\
\midrule
  \textbf{Macro F1} & 0.1813 & \textbf{0.1813} & +0.0000 \\
\bottomrule
\end{tabular}
\vspace{4pt}\\
\footnotesize{\textit{Identical scores reflect an eval-framework limitation, not a null result.
Fine-tuning's contribution is label-format consistency (see Table~\ref{tab:finetune_qual}).}}
\end{table}

% Qualitative effect of QLoRA fine-tuning — label consistency evidence
% Card: test_germany_0015.jpg
\begin{table}[h]\centering
\caption{Qualitative Effect of QLoRA Fine-tuning on VLM Output Format}
\label{tab:finetune_qual}
\begin{tabular}{lll}
\toprule
Field & Base Model Output & Fine-tuned Output \\
\midrule
Name      & \texttt{Name | Leonie D\"{o}rr}         & \texttt{Name | Leonie D\"{o}rr} \\
DOB       & \texttt{Date of Birth | 01.08.1991}    & \texttt{Date of Birth | 01.08.1991} \\
ID Number & \texttt{\textit{[empty]} | M898585166} & \texttt{ID No. | M898585166} \\
Signature & \texttt{Signature | [none]}             & \texttt{Signature | [none]} \\
\bottomrule
\end{tabular}
\vspace{2pt}\\
\footnotesize{The base model detects PII values correctly but omits field labels on
structured fields (row~3). Fine-tuning teaches consistent \texttt{LABEL\,|\,VALUE}
output format required by the downstream pipeline: an empty label is normalised to
\texttt{other} by \texttt{normalize\_cat()} and scored as a false positive regardless
of whether the value is correct. Quantitative F1 appears identical between the two
models (Table~\ref{tab:finetune}) because the evaluation framework matches on values,
not on the label\,$\rightarrow$\,category mapping the live pipeline depends on.}
\end{table}


Table~\ref{tab:finetune_qual} shows a representative
comparison for a single German identity card.
The base model correctly detects the ID number value
(\texttt{M898585166}) but outputs an \emph{empty
label} for that field --- yielding a raw prediction
of \texttt{[empty] | M898585166}.
The fine-tuned model outputs \texttt{ID No. |
M898585166} --- the correctly labelled form required
by the downstream BBRefinerAgent and ContextAgent.

This distinction matters because the evaluation
framework maps predictions to PII categories via
the label field: an empty label is classified as
\texttt{other} and counted as a false positive,
while the correct label is mapped to
\texttt{id\_numbers} and matched against the ground
truth.
Two predictions with identical values but different
labels produce different F1 contributions.

The fine-tuning contribution is therefore
\textbf{output format consistency}: the adapter
trains the model to emit consistently parseable
\texttt{LABEL\,|\,VALUE\,|\,BBOX} triples for all
detected fields, not just the most visually salient
ones.
This is a necessary condition for reliable pipeline
operation --- a model that intermittently omits field
labels would require additional heuristic recovery
logic in every downstream agent.

We note two additional qualitative improvements
observed post fine-tuning.
First, \emph{face detection generalises}: the base
model rarely labels the face photograph region
explicitly, while the fine-tuned model reliably
identifies it as a Face entity with a reasonable
bounding box.
Since face photographs are not labelled as explicit
PII in the Claude teacher annotations (the teacher
labels text fields only), this generalisation appears
to be an emergent capability from learning the
structured field-labelling task.
Second, \emph{label vocabulary stabilises}: the base
model uses inconsistent label strings (\texttt{Full
Name}, \texttt{Surname}, \texttt{Given Name},
\texttt{Name}) for the same field type across
different cards, requiring a broad normalisation
dictionary in the ContextAgent.
The fine-tuned model converges to the label
conventions present in the training data,
substantially reducing the normalisation burden.

% ─────────────────────────────────────────────────────────
\subsection{Synthetic-to-Real Transfer}
% ─────────────────────────────────────────────────────────

A standard concern with synthetic training data is
distribution shift: models trained on clean synthetic
images may fail on the photographic noise, printing
artefacts, and typographic diversity of real documents.

We address this directly in Section~\ref{sec:evaluation}
by evaluating Omni-Shield on 19 real specimen identity
documents (Section~\ref{sec:real_eval}).
The finding is that Macro F1 on real documents
(0.269) \emph{exceeds} performance on synthetic test
cards (0.187).
We attribute this counterintuitive result to two
factors: the synthetic test set uses greater layout
variation than the training set (harder than real
specimen cards in our evaluation), and the visual
features that the VLM uses for layout understanding
--- field label typography, spatial proximity of
label-value pairs, card background colour --- are
faithfully reproduced by the Pillow renderer.

The largest performance gap between synthetic and
real evaluation is in the Names category
(F1: 0.312 synthetic, 0.182 real), which we
attribute to typographic diversity in real documents
--- ALL CAPS names, diacritics, and multi-word
surnames are underrepresented in the Faker-generated
training distribution.
Expanding training coverage of name formats is
identified as the highest-priority data augmentation
for future work.
